%--IPFW ECE405 Template--%
\documentclass{ipfwpaper}
\svnid{$Id$}
\usepackage{cite}

%******Cover Page*******%
%Step 1: Add your names and emails
%Step 2: Add your advisors names and emails
%Step 3: Add your title


%Step 1
%--Authors Names and Emails--%
\author{Author 1\formatemail{Email 1} Author 2\formatemail{Email 2} Author 3\formatemail{Email 3}}  

%Step 2
%--Advisors Names and Emails--%
\advisor{Advisor 1\formatemail{Email 1} Advisor 2\formatemail{Email 2}}

%Step 3
%--Title of project--%
\title{CMPEN 480/ EE 400:Senior Design Report} 

\courseadvisor { Dr. Abdallah S. Abdallah \\
Assistant Professor of Electrical \& Computer Engineering}


%---Acronyms---%
%Use the acronyms.tex file to keep track of your acronyms. A list  will be automatically generated in the appendix. See the Acknowledgements section for an example.
\loadglsentries{acronyms}
\makeglossaries
%-----------------%



%-----Beginning of Document----%
\begin{document} 
\maketitle %Makes the cover page




%---Revision Tracker---%
\startTable
\AddRevision{Number}{Date}{Add explanation}
\AddRevision{}{}{}
\stopTable


%-----Table of Contents-----%
\tableofcontents
\listoftables
\listoffigures
\printglossary[title=List of Acronyms ,type=\acronymtype]
\clearpage
%--end:Table of Contents--%



\section*{Acknowledgments}
\addcontentsline{toc}{section}{Acknowledgments}%adds acknowledgements section to the table of contents


Here is a sample for how to use acronyms. "Defining a \gls{fr} and a \gls{ps} provides a language for system design. This language is an important aspect of \gls{csd}."


\input{Abstract}


%%%%%%  PROBLEM STATEMENT  %%%%%%                                               
\input{ProblemStatement}
%%%%%%  Conceptual Design %%%%%%  
%% Report (1) only includes Problem Statements, so all other sections are commented except references Section
%\clearpage
%\input{ConceptualDesign}

%\clearpage
%%%%%%   Evaluation of Design Alternatives %%%%%%  
%\input{EvaluationofDesignAlternatives}
%%%%%%   Detailed Design %%%%%%  

%\clearpage
%\input{DetailedDesign}
%%%%%% Cost %%%%%% 
%\clearpage
%\input{Cost}
%%%%%%  %%%%%%  Test Development and Execution %%%%%%  %%%%%%  
%\clearpage
%\input{RiskManagement}
%%%%%%   Project Timeline  %%%%%%  
%\clearpage
%\input{ProjectPlan}

        
%%%%%%   Conclusions  %%%%%%  
%\clearpage
%\input{ProjectPlan}





%----References-----%
\clearpage

\addcontentsline{toc}{section}{References}
%adds reference section to the table of contents
%Please refer to the OverLeaf Primer to see your options for how to keep track of references.
Use the OverLeaf Primer to understand the different ways to format the references. 
\bibliographystyle{IEEEtran} 
%Formats the bibliography to meet the IEEE standards
\bibliography{MyReferences.bib}
%This command calls a file named "your_references.bib". You will upload this file, and can rename as necessary. 


\end{document}
